\TieudeT{PHONG TỎA VẬT LÍ 11 - DAO ĐỘNG VÀ SÓNG}
\subsubsection*{PHẦN I. Thí sinh trả lời câu hỏi từ câu 1 đến câu 18. Mỗi câu hỏi thí sinh chỉ chọn một phương án}
\setcounter{ex}{0}
	\begin{ex} %[3P1Y1-1][Loai1] 
		Một chất điểm dao động điều hòa với phương trình $x=A\cos\left(\omega t+\varphi\right)$, trong đó $\omega $ có giá trị dương. Đại lượng $\omega$ gọi là
		\choice
		{biên độ dao động}
		{chu kì của dao động}
		{\True tần số góc của dao động}
		{pha ban đầu của dao động}
		\loigiai{
			$\omega$ được gọi là tần số góc của dao động.
		}
		%\haidong{5}
	\end{ex}
	\begin{ex}
		Một chất điểm dao động điều hoà có phương trình li độ theo thời gian là $x=6\cos(4\pi t+\dfrac{\pi}{3})\ (cm)$. Chu kì của dao động là
		\choice
		{$4\ s$}
		{$2\ s$}
		{$0,25\ s$}
		{\True $0,5\ s$}
		\loigiai{
			\begin{itemize}
				\item Chu kì của dao động bằng: $T=\dfrac{2\pi}{\omega}=0,5\ s$
			\end{itemize}
		}	
		%\haidong{5}
	\end{ex}
	\begin{ex} %[3P1Y2-1][Loai1] 
		Một vật dao động điều hòa trên trục Ox. Vận tốc của vật
		\choice
		{luôn có giá trị không đổi}
		{luôn có giá trị dương}
		{là hàm bậc hai của thời gian}
		{\True biến thiên điều hòa theo thời gian}
		\loigiai{
			Vận tốc của vật biến thiên điều hòa theo thời gian.
		}	
		%\haidong{5}
	\end{ex}
	%\loai{Cho li độ x tìm v}
	\begin{ex} %[3P1B2-2][Loai1] 
		Vật dao động điều hoà với biên độ $A = 5\ cm$, tần số $f = 4\ Hz$. Tốc độ vật khi có li độ $x = 3\ cm$ là
		\choice 
		{$\left|v\right| = 2\pi\ cm/s$}
		{$\left|v\right|= 16\pi\ cm/s$}
		{\True $\left|v\right|= 32\pi\ cm/s$}
		{$\left|v\right|= 64\pi\ cm/s$}
		\loigiai{
			\begin{itemize}
				\item Ta có: $\omega = 2\pi f = 8\pi\ rad/s.$
				\item Biểu thức độc lập: $A^2 = x^2 + \dfrac{v^2}{\omega^2}\Rightarrow |v| = \omega\sqrt{A^2-x^2} = 8\pi\sqrt{0,05^2-0,03^2} = 0,32\pi\ m/s = 32\pi\ cm/s$.
			\end{itemize}
		}
		%\haidong{10}
	\end{ex}
	%\loai{Cho pha, tìm li độ và chiều}
	\begin{ex}%[3P1B4-1][Loai1]
		Một vật nhỏ dao động điều hòa với biên độ A trên trục Ox. Khi pha của vận tốc là $0$ thì vật
		\choice
		{ở biên dương $x = A$}
		{đi qua vị trí cân bằng theo chiều âm}
		{\True đi qua vị trí cân bằng theo chiều dương}
		{ở biên âm $x = -A$}
		\loigiai{
			Ta có: $\phi _x=\phi _v-\dfrac{\pi}{2} = - \dfrac{\pi}{2} \Rightarrow x = 0\quad (+)$.
		}
		%\haidong{8}
	\end{ex}
	%\loai{Khoảng thời gian ngắn nhất từ li độ $x_1$ đến li độ $x_2$ vị trí bất kì}
	\begin{ex}%[3P1K8-1][Loai1]
		Một vật dao động điều hòa với chu kì $2\ s$, biên độ $A$. Khoảng thời gian ngắn nhất vật đi từ vị trí $0,6A$ đến vị trí $-0,8A$ là
		\choice
		{$0,41\ s$}
		{$0,59\ s$}
		{$0,205\ s$}
		{\True $0,5\ s$}
		\loigiai{
			\begin{itemize}
				\item Sử dụng công thức học về khoảng thời gian vật dao động giữa vị trí cân bằng và li độ x không đặc biệt là:
				$\dfrac{\arcsin \left( \dfrac{\left| x \right|}{A} \right)}{\omega }=\dfrac{T.\arcsin \left( \dfrac{\left| x \right|}{A} \right)}{2\pi}$.\\
				$\Rightarrow$ Thời gian cần tìm là :  $\dfrac{T.\arcsin \left( 0,6 \right)}{2\pi}+\dfrac{T.\arcsin \left( 0,8 \right)}{2\pi}=\dfrac{T}{2\pi}.\dfrac{\pi}{2}=\dfrac{T}{4}$. 
			\end{itemize}
		}
		
		%\haidong{15}
	\end{ex}
	%\loai{Khoảng thời gian vật cách vị trí cân bằng nhỏ hơn b trong một chu kì}
	
	\begin{ex} %[3P1B8-3][Loai1]  
		Cho một chất điểm đang dao động điều hòa với chu kì bằng $2\ s$ và biên độ bằng $6\ cm$. Trong một chu kì dao động, tổng thời gian mà khoảng cách từ chất điểm tới vị trí cân bằng nhỏ hơn $3\sqrt{3}\ cm$ là
		\choice
		{\True $\dfrac{4}{3}\ s$}
		{$\dfrac{2}{3}\ s$}
		{$1\ s$}
		{$\dfrac{3}{4}\ s$}
		\loigiai{
			\immini{\begin{itemize}
					\item 
					Ta có: $|x|< +3\sqrt{3} \Rightarrow -3\sqrt{3}\ cm< x< +3\sqrt{3}\ cm$.
					\item Ta biểu diễn các điểm có li độ $3\sqrt{3}$ trên đường tròn lượng giác như hình vẽ.
					\item Trong một chu kì dao động, khoảng cách từ chất điểm tới vị trí cân bằng nhỏ hơn $3\sqrt{3}$ khi nó đi từ vị trí $M_1$ đến $M_2$ và từ $M_3$ đến $M_0\Rightarrow t = \dfrac{T}{3}+\dfrac{T}{3} = \dfrac{2T}{3} = \dfrac{4}{3}\ s$.
			\end{itemize}}{\begin{tikzpicture}[scale=1,thick,>=stealth']
					\tkzInit[ymin=-3,ymax=3,xmin=-3,xmax=3]\tkzClip
					\tkzDefPoints{-2.5/0/m, 2.5/0/n, 0/0/o, 0/2.5/p, 0/-2.5/q}
					\tkzDefShiftPoint[o](-150:2){0}
					\tkzDefShiftPoint[o](-30:2){1}
					\tkzDefShiftPoint[o](30:2){2}
					\tkzDefShiftPoint[o](150:2){3}
					\tkzDefPointBy[projection=onto m--n](1) \tkzGetPoint{h}
					\tkzDefPointBy[projection=onto m--n](0) \tkzGetPoint{h'}
					\draw(o)circle(2cm);
					\tkzDrawSegments[dashed](0,3 1,2)
					\tkzDrawSegments(p,q)
					\tkzDrawSegments[->](m,n)
					\tkzDrawSegments[blue,->](o,1 o,2 o,0 o,3)
					\tkzDrawPoints[size=3](o,1,h,0,2,3,h')
					\tkzMarkAngles[size=0.7cm,arrows=->](2,o,3)
					\tkzMarkAngles[size=0.7cm,arrows=->](0,o,1)
					\node at (0)[below left]{$M_3$};
					\node at (1)[below right]{$M_0$};
					\node at (2)[above right]{$M_1$};
					\node at (3)[above left]{$M_2$};
					\node at (2,0)[above right]{$6$};
					\node at (h')[below]{$-3\sqrt{3}$};
					\node at (h)[below]{$3\sqrt{3}$};
			\end{tikzpicture}}
		}%<MyLT1>
		%\haidong{15}
	\end{ex}
	%\loai{Thời gian tốc độ thỏa mãn giá trị cho trước trong một chu kì}
	\begin{ex}%[3P1K9-1][Loai1]
		Cho một chất điểm dao động điều hòa với biên độ bằng $4\ cm$, chu kì bằng $0,5\ s$. Trong mỗi chu kì, khoảng thời gian mà tốc độ chuyển động của vật lớn hơn $8\pi\sqrt{2}\ cm/s$ là
		\choice
		{$\dfrac{3}{4}\ s$}
		{$\dfrac{1}{8}\ s$}
		{$\dfrac{1}{6}\ s$}
		{\True $\dfrac{1}{4}\ s$}
		\loigiai{
			\immini{\begin{itemize}
					\item Ta có: $\omega = \dfrac{2\pi}{0,5} = 4\pi \Rightarrow v_{\max} = 4\pi. 4 = 16\pi\ cm/s$.
					\item Theo đề: $|v|>8\pi\sqrt{2}\ cm/s \Rightarrow v>8\pi\sqrt{2}\ cm/s$ hoặc $v< -8\pi\sqrt{2}\ cm/s$.
					\item Trong 1 chu kì, vật chuyển động với tốc độ lớn hơn $8\pi\sqrt{2}$ khi nó đi từ điểm $M_0$ đến $M_1$ và từ điểm $M_3$ đến $M_4$ trên đường tròn lượng giác như hình vẽ\\
					\\
					$\Rightarrow \Delta \varphi = \dfrac{\pi}{2}+\dfrac{\pi}{2} = \pi\Rightarrow t = \dfrac{\Delta \varphi}{\omega} = \dfrac{1}{4}\ s$.
			\end{itemize}}{\begin{tikzpicture}[scale=1,thick,>=stealth']
					\tkzInit[ymin=-3,ymax=3,xmin=-3,xmax=3]\tkzClip
					\tkzDefPoints{-2.5/0/m, 2.5/0/n, 0/0/o, 0/2.5/p, 0/-2.5/q}
					\tkzDefShiftPoint[o](-135:2){1}
					\tkzDefShiftPoint[o](-45:2){2}
					\tkzDefShiftPoint[o](45:2){3}
					\tkzDefShiftPoint[o](135:2){4}
					\tkzDefPointBy[projection=onto m--n](2) \tkzGetPoint{h}
					\tkzDefPointBy[projection=onto m--n](1) \tkzGetPoint{h'}
					\draw(o)circle(2cm);
					\tkzDrawSegments[dashed](1,4 2,3)
					\tkzDrawSegments(p,q)
					\tkzDrawSegments[->](m,n)
					\tkzDrawSegments[blue,->](o,1 o,2 o,4 o,3)
					\tkzDrawPoints[size=3](o,1,h,4,2,3,h')
					\tkzMarkAngles[size=0.7cm,arrows=->](4,o,1 2,o,3)
					\node at (1)[below left]{$M_2$};
					\node at (2)[below right]{$M_3$};
					\node at (3)[above right]{$M_4$};
					\node at (4)[above left]{$M_1$};
					\node at (2.5,0)[below right]{$v$};
					\node at (2,0)[above right]{$16\pi$};
					\node at (h')[below]{$-8\pi\sqrt{2}$};
					\node at (h)[below]{$8\pi\sqrt{2}$};
		\end{tikzpicture}}}
		
		
		%\haidong{15}
	\end{ex}
	%\loai{Bài toán tổng hợp}
	\begin{ex}%[3P1G9-1][Loai1]
		Một vật nhỏ dao động điều hòa theo phương ngang với tần số góc $\omega$. Tại thời điểm $t = 0$, vật nhỏ qua vị trí cân bằng theo chiều dương. Tại thời điểm $t = 1,475\ s$, vận tốc $v$ và li độ $x$ của vật nhỏ thỏa mãn $v=\omega \left| x \right|\sqrt{3}$ lần thứ $10$. Tần số của vật xấp xỉ
		\choice
		{$3,1\ Hz$}
		{$\dfrac{10\pi}{3}\ Hz$}
		{\True $3,3\ Hz$}
		{$20,7\ Hz$}
		\loigiai{
			\begin{itemize}
				\item Khi $v=\omega \left| x \right|\sqrt{3}$ (v >0) $\Rightarrow {{x }^2}+\dfrac{v^2}{\omega^2}=A^2\to x=\pm \dfrac{A}{2}\ (+)$.
				\item Thấy mỗi chu kì vật qua $x=\pm \dfrac{A}{2}(+)$ 2 lần, do đó, tách $10 = 8 + 2$.
				\item Vậy sau 4T vật qua $x=\pm \dfrac{A}{2}(+)$ 8 lần và quay lại trạng thái tại $t = 0$: $x = O(+)$, vật đi qua 2 lần nữa theo diễn biến trục thời gian bên dưới mất $\dfrac{T}{4}+\dfrac{T}{2}+\dfrac{T}{6}$.
				\begin{center}
					\begin{tikzpicture}
						\tkzDefPoints{0/0/o, 5/0/a, 4.325/0/b, 3.525/0/c, 2.5/0/d,  -5/0/h, -4.325/0/g, -3.525/0/f, -2.5/0/e, -6/0/x', 6/0/x}
						\tkzDefShiftPoint[a](-90:0.5cm){a1}
						\tkzDefShiftPoint[a](-90:1cm){a2}
						\tkzDefShiftPoint[o](-90:0.5cm){o1}
						\tkzDefShiftPoint[h](-90:1cm){h2}
						\tkzDefShiftPoint[h](-90:1.5cm){h3}
						\tkzDefShiftPoint[e](-90:1.5cm){e3}
						\tkzDrawSegments[blue,->,very thick](x',x)
						\tkzDrawPoints[size=5,fill=black](o,a,b,c,d,e,f,g,h)
						\tkzLabelPoint[above](o){\tiny $O$}
						\tkzLabelPoint[above](a){\tiny$A$}
						\tkzLabelPoint[above](b){\tiny$\dfrac{A\sqrt{3}}{2}$}
						\tkzLabelPoint[above](c){\tiny$\dfrac{A\sqrt{2}}{2}$}
						\tkzLabelPoint[above](d){\tiny$\dfrac{A}{2}$}
						\tkzLabelPoint[above](h){\tiny$-A$}
						\tkzLabelPoint[above](g){\tiny $-\dfrac{A\sqrt{3}}{2}$}
						\tkzLabelPoint[above](f){\tiny$-\dfrac{A\sqrt{2}}{2}$}
						\tkzLabelPoint[above](e){\tiny$-\dfrac{A}{2}$}
						\tkzLabelPoint[above](x){\tiny$x$}
						\tkzDrawSegments[->,very thick,red](o1,a1 a2,h2 h3,e3)
						\tkzDrawSegments[very thick,red](a1,a2 h2,h3)
					\end{tikzpicture}
				\end{center}
				\item Thời điểm $t = 1,475\ s = 4T +\dfrac{T}{4}+\dfrac{T}{2}+\dfrac{T}{6} \Leftrightarrow T = 0,3\ s \Rightarrow \omega =\dfrac{20\pi}{3} \Rightarrow f=\dfrac{10}{3}\ Hz$.
			\end{itemize}
		}
		%\haidong{15}
	\end{ex}
	%\loai{Tính quãng đường vật đi được trong quãng thời gian $\Delta t=(2n+1)\dfrac{T}{2}$}
	\begin{ex}%[3P1BD-1][Loai1] 
		Một chất điểm dao động điều hoà với tần số góc $\omega=20\ rad/s$ và biên độ $A = 6\ cm$. Chọn gốc thời gian lúc vật đi qua vị trí cân bằng. Quãng đường vật đi được trong $0,05\pi\ s$ là
		\choice
		{$24\ cm$}
		{$9\ cm$}
		{$6\ cm$}
		{\True $12\ cm$}
		\loigiai{
			\begin{itemize}
				\item Trong $0,05\pi = \dfrac{T}{2}$ Vật đi được quãng đường $2A = 12\ cm$ dù cho gốc thời gian là ở lúc nào.
				\item 
			\end{itemize}
		}
		%\haidong{20}
	\end{ex}
	%\loai{Tính chu kì, tần số}
	\begin{ex}%[3P1B3-2]
		Một con lắc lò xo, vật nặng có khối lượng $m = 250\ g$, lò xo có độ cứng $k = 100\ N/m$. Tần số dao động của con lắc xấp xỉ
		\choice
		{$20\ Hz$}
		{\True $3,18\ Hz$}
		{$6,28\ Hz$}
		{$5\ Hz$}
		\loigiai{
			Tần số dao động của con lắc lò xo là: $f = \dfrac{1}{{2\pi}}\sqrt{\dfrac{k}{m}} = \dfrac{1}{{2\pi}}\sqrt {\dfrac{100}{0,25}} = 3,18\ Hz$.
		}
		%\haidong{5}
	\end{ex}
	%\loai{Loại khác}
	\begin{ex}%[3P1-19.1K][Loai1]
		
		\immini[thm]{Sự phụ thuộc của li độ theo thời gian của một chất điểm dao động điều hòa được biểu diễn như hình vẽ. Tại điểm nào trong các điểm M, N, K, H có vectơ gia tốc và vectơ vận tốc ngược chiều nhau?
			
			\choice
			{Điểm M}
			{\True Điểm K}
			{Điểm N}
			{Điểm H}}{\begin{tikzpicture}[draw=Mapcolor,scale=1,thick,>=stealth']
				\draw[->] (0,0)--(7,0)node[below]{$t\ (s)$};
				\draw[->] (0,-2.3)--(0,2.3)node[right]{$x\ (cm)$};
				\draw[dashed] (0,1.5)--(6,1.5)--(6,0)(0,-1.5)--(6,-1.5)--(6,0);
				\draw [Mapcolor, line width = 1.2pt, domain=0:6, samples=500]%
				plot (\x, {1.5*cos(\x*120)});
				\draw(0.75,0)circle(1pt)node[above right]{M};
				\draw(1.5,-1.5)circle(1pt)node[below]{N};
				\draw(2.5,0.75)circle(1pt)node[above left]{K};
				\draw(3.5,0.75)circle(1pt)node[above right]{H};
		\end{tikzpicture}}
		\loigiai{
			Từ đồ thị ta xét các vị trí:
			\begin{itemize}
				\item Tại M vật đang đi qua vị trí cân bằng theo chiều âm $\Rightarrow$ gia tốc của vật có độ lớn $ = 0$.
				\item Tại N vật đang ở vị trí biên âm $\Rightarrow$ vận tốc của vật có độ lớn $=0$.
				\item Tại K vật đang có li độ x và đang chuyển động từ vị trí cân bằng ra biên dương\\ $\Rightarrow$ Tại K vectơ gia tốc hướng ngược chiều dương, vectơ vận tốc hướng cùng chiều dương
				\item Tại H vật đang có li độ x và đang chuyển động từ vị trí biên dương về vị trí cân bằng \\
				$\Rightarrow$ Tại K vectơ gia tốc hướng ngược chiều dương, vectơ vận tốc hướng ngược chiều dương.
			\end{itemize}
		}
		%\haidong{20}
	\end{ex}
	\begin{ex} %[3P1-20.1G][Loai1]
		\immini[thm]{Một chất điểm dao động điều hòa có pha dao động của li độ quan hệ với thời gian được biểu diễn như hình vẽ. Quãng đường chất điểm đi được từ thời điểm $t_3$ đến thời điểm $t_4$ là $10\ cm$ và $t_2 - t_1 = 0,5\ s$. Gia tốc của chất điểm tại thời điểm $t = 2018\ s$ \textbf{gần nhất} với giá trị nào sau đây?
			\haicot
			{$17\ cm/s^2$}
			{\True $22\ cm/s^2$}
			{$20\ cm/s^2$}
			{$14\ cm/s^2$}}{\begin{tikzpicture}[draw=Mapcolor,scale=1,thick,>=stealth']
				\draw[->] (0,0)--(4.5,0)node[above]{$t\ (s)$};
				\draw[->] (0,-1.3)--(0,2.5)node[right]{$\omega t+\varphi\ (rad)$};
				\foreach \x in {0.5,1,...,4}
				{\draw [dashed,thin,Mapcolor] (\x,-1) --(\x,2);}
				\foreach \y in {-1,-0.5,0.5,1,...,2}
				{\draw [dashed,thin,Mapcolor] (0,\y) --(4,\y);}
				\node at (0,0)[left]{$O$};
				\node at (0,0.5)[left]{$\dfrac{\pi}{3}$};
				\draw[fill=white] (0.5,-1)circle(1pt)node[below]{$t_1$};
				\draw[fill=white] (1,-1)circle(1pt)node[below]{$t_2$};
				\draw[fill=white] (2,-1)circle(1pt)node[below]{$t_3$};
				\draw[fill=white] (3,-1)circle(1pt)node[below]{$t_4$};
				\draw[fill=white] (0,0.5)circle(1pt);
				\draw[fill=white] (0,1)circle(1pt);
				\draw[fill=white] (0,1.5)circle(1pt);
				\draw[fill=white] (0,2)circle(1pt);
				\draw[fill=white] (0,0)circle(1pt);
				\draw[Mapcolor,very thick](0,-1)--(3,2);
				\draw[fill=white] (0,-1)circle(1pt);
				\draw[fill=white] (0.5,-0.5)circle(1pt);
				\draw[fill=white] (1,0)circle(1pt);
				\draw[fill=white] (2,1)circle(1pt);
				\draw[fill=white] (3,2)circle(1pt);
			\end{tikzpicture}
		}
		\loigiai{
			\begin{itemize}
				\item Ta có $t_2 - t_1 = 0,5\ s\ \Rightarrow$ trên Ot mỗi khoảng tương ứng 0,5 s.
				\item Khi $t = 0$ thì $(\omega.0 + \varphi) = - \dfrac{2\pi}{3} \Rightarrow \varphi = - \dfrac{2\pi}{3}$.
				\item Khi $t = t_2 = 1\ s$ thì $(\omega.1 + \varphi) = 0 \Rightarrow \omega = \dfrac{2\pi}{3}\ rad/s \Rightarrow T = 3\ s$.
			\end{itemize}
			\immini{\begin{itemize}
					\item Phương trình dao động có dạng $x = A\cos(\dfrac{2\pi}{3}t - \dfrac{2\pi}{3}) cm$
					\item Khi $t = t_3 = 2\ s$ thì pha là $\dfrac{2\pi}{3}\ rad$.
					\item Khi $t = t_4 = 3\ s$ thì pha là $\dfrac{4\pi}{3}\ rad \Rightarrow$ góc quét $\Delta\varphi = \dfrac{2\pi}{3}\ rad$ trên vòng tròn.
					\item Dễ dàng tính được $S = A = 10\ cm$.
					$\Rightarrow$ Phương trình dao động $x = 10\cos(\dfrac{2\pi}{3}t - \dfrac{2\pi}{3}) cm$
					\item Gia tốc $a = - \left(\dfrac{2\pi}{3}\right)^2.10\cos(\dfrac{2\pi}{3}t-\dfrac{2\pi}{3})\ cm/s^2$. 
					\item Khi $t = 2018 s$ thì $a = 21,9\ cm/s^2$.
			\end{itemize}}{\begin{tikzpicture}[scale=1,thick,>=stealth']
					\tkzInit[ymin=-3,ymax=3,xmin=-3,xmax=3]\tkzClip
					\tkzDefPoints{-2.5/0/m, 2.5/0/n, 0/0/o, 0/2.5/p, 0/-2.5/q}
					\tkzDefShiftPoint[o](240:2){2}
					\tkzDefShiftPoint[o](120:2){1}
					\tkzDefShiftPoint[o](90:2){0}
					\tkzDefPointBy[projection=onto m--n](1) \tkzGetPoint{h}
					\draw(o)circle(2cm);
					\tkzDrawSegments[dashed](1,2)
					\tkzDrawSegments(p,q)
					\tkzDrawSegments[->](m,n)
					\tkzDrawSegments[blue,->](o,1 o,2)
					\tkzDrawPoints[size=3](o,1,h,2)
					\node at (1)[above left]{$t_3$};
					\node at (2)[below left]{$t_4$};
					\node at (2,0)[below right]{$A$};
					\node at (h)[below left]{$-\dfrac{A}{2}$};
				\end{tikzpicture}
			}
		}
		%\haidong{20}
	\end{ex}
	
	\begin{ex}
		Một con lắc lò xo gồm vật nhỏ và lò xo nhẹ có độ cứng k, đang dao động điều hòa theo phương nằm ngang. Mốc thế năng ở vị trí cân bằng. Khi vật qua vị trí có li độ x thì thế năng của con lắc có biểu thức nào sau đây?
		\choice
		{\True $W_t=\dfrac{1}{2}kx^2$}
		{$W_t=kx$}
		{$W_t=\dfrac{1}{2}kx$}
		{$W_t=kx^2$}
		%\haidong{5}
	\end{ex}
	\begin{ex}%[3P1BF-1][Loai1]
		Tại nơi có gia tốc trọng trường g, một con lắc đơn có sợi dây dài $\ell $ đang dao động điều hoà. Tần số dao động của con lắc là
		\choice
		{$2\pi \sqrt{\dfrac{\ell}{g}}$}
		{$2\pi \sqrt{\dfrac{g}{\ell}}$}
		{\True $\dfrac{1}{2\pi}\sqrt{\dfrac{g}{\ell}}$}
		{$\dfrac{1}{2\pi}\sqrt{\dfrac{\ell}{g}}$}
		\loigiai{
			Tần số dao động của con lắc đơn là: $f=\dfrac{1}{2\pi}\sqrt{\dfrac{g}{\ell}}$.
		}
		%\haidong{5}
	\end{ex}
	%\loai{Tìm chu kì}
	
	\begin{ex}%[3P1K6-2][Loai1]%Bài 15 dạng 1
		Một con lắc lò xo treo thẳng đứng, dao động điều hòa dọc theo quỹ đạo dài $6\ cm$. Khi vật ở vị trí cao nhất, lò xo bị giãn $1\ cm$. Lấy $g = 10\ m/s^2,\ \pi^2 = 10$. Chu kì dao động của con lắc bằng
		\choice
		{$0,5\ s$}
		{$0,6\ s$}
		{\True $0,4\ s$}
		{$0,3\ s$}
		\loigiai{
			\begin{itemize}
				\item Ở vị trí cao nhất (biên bên trên), lò xo bị giãn 1 cm
				\item Vật ở dưới vị trí lò xo tự nhiên 1 cm $\Rightarrow \Delta \ell_0 = A + 1\  cm = 4\ cm \Rightarrow T = 0,4\ s$.	
			\end{itemize}
		}
		%\haidong{10}
	\end{ex}
	\begin{ex} %[3P1-16.1Y][Loai1]   
		Một con lắc đơn gồm dây treo dài $\ell$ và vật có khối lượng là $m$. Con lắc treo tại nơi có gia tốc rơi tự do là g. Kích thích con lắc dao động với biên độ góc $\alpha _0$ rất nhỏ. Bỏ qua mọi ma sát và lực cản. Biểu thức năng lượng dao động của con lắc là
		\choice
		{$2mg\ell\alpha _0^2$}
		{\True $\dfrac{1}{2}{mg\ell}\alpha _0^2$}
		{$mg\ell\alpha _0^2$}
		{$\dfrac{2mg}{\ell}\alpha _0^2$}
		\loigiai{
			Cơ năng của con lắc được xác định bằng biểu thức: $W=\dfrac{1}{2}{mg\ell}\alpha _0^2$.
		}
		%\haidong{5}
	\end{ex}
	\begin{ex} %[3P1-16.3][Loai1]  
		Một con lắc đơn dao động điều hòa với biên độ góc $0,1\ rad$ ở một nơi có gia tốc trọng trường $g = 10\ m/s^2$. Vào thời điểm ban đầu vật đi qua vị trí có li độ dài $8\ cm$ và có vận tốc $20\sqrt{3}\ cm/s$. Tốc độ cực đại của vật dao động là
		\choice
		{$0,8\ m/s$}
		{$0,2\ m/s$}
		{\True $0,4\ m/s$}
		{$1\ m/s$}
		\loigiai{
			Áp dụng công thức độc lập thời gian, ta có:
			\begin{align*}
				S_0^2=s^2+\dfrac{v^2}{\omega ^2}\Leftrightarrow &\left(\ell \alpha_0\right)^2=s^2+\dfrac{\ell.v^2}{g}\\
				\Leftrightarrow &\left(\ell.0,1\right)^2=0,08^2+\dfrac{\ell.0,04.3}{10}\Rightarrow \ell=1,6\ m
				\Rightarrow v_{\max}=\omega S_0=\sqrt{\dfrac{g}{\ell}}.\ell \alpha_0=0,4\ m/s.
			\end{align*}1
		}
		%\haidong{10}
	\end{ex}
	

\subsubsection*{PHẦN 2. Thí sinh trả lời từ câu 1 đến câu 4. Trong mỗi ý a), b), c), d) ở mỗi câu, thí sinh chọn đúng hoặc sai.}
\setcounter{ex}{0}
\begin{ex}
	Một vật dao động điều hòa có phương trình $x=4 \cos (4 \pi t+\dfrac{\pi}{6})\ cm$.
	\choiceTF[t]
	{\True Biên độ của vật là $4 \ cm$}
	{Từ vị trí cân bằng ra biên thì vật chuyển động nhanh dần đều}
	{Thời điểm lần thứ $3$ khi vật qua vị trí có li độ $2 \ cm$ theo chiều dương là $\dfrac{9}{8} \ s$}
	{Các thời điểm vật chuyển động qua vị trí có toạ độ $x=-2 \ cm$ theo chiều dương $O x$ là $t=0,5+2 k\ (s)\ (k=1,2,3, \ldots)$} \loigiai{
		\begin{itemchoice}
			\itemch Đúng.
			\itemch Sai.
			\itemch Sai.
			\itemch Sai.
		\end{\itemchoice}
	}
	%\haidong{25}
\end{ex}
\begin{ex}
	Một vật đang dao động điều hòa theo phương trình $x=-5 \cos 5 \pi t\ cm$.
	\choiceTF[t]
	{Pha ban đầu của li độ bằng $0\ rad$}
	{\True Biên độ dao động là $5 \ cm$}
	{\True Tần số dao động là $2,5 \ Hz$}
	{\True Chu kì dao động là $0,4 \ s$}
	\loigiai{
		\begin{itemchoice}
			\itemch Sai.
			\itemch Đúng.
			\itemch Đúng.
			\itemch Đúng.
		\end{\itemchoice}
	}
	%\haidong{25}
\end{ex}

\begin{ex}
	Trong dao động điều hoà cơ năng luôn được bảo toàn và có sự chuyển hoá qua lại giữa động năng và thế năng.
	\choiceTF[t]
	{\True Cơ năng của vật dao động điều hoà bằng tổng động năng và thế năng ở thời điểm bất kì}
	{Cơ năng của vật dao động điều hoà biến thiên cùng tần số với động năng}
	{Cơ năng của vật dao động điều hoà tăng hai lần khi biên độ tăng hai lần}
	{\True Cơ năng của vật dao động điều hoà bằng động năng khi vật ở vị trí cân bằng}
	\loigiai{
		\begin{itemchoice}
			\itemch Đúng.
			\itemch Sai.
			\itemch Sai.
			\itemch Đúng.
		\end{\itemchoice}
	}
	%\haidong{23}
\end{ex}
\begin{ex}
	Một vật có khối lượng $m=100 \ g$ được gắn vào lò xo có độ cứng $k=100 \ N/m$. Kích thích cho con lắc dao động điều hòa theo phương ngang. Trong quá trình dao động thì chiều dài lớn nhất của lò xo là $5 \ cm$.
	\choiceTF[t]
	{Khi vật đi đến biên thì vận tốc vật lớn nhất}
	{Tần số góc dao động là $5 \ Hz$}
	{Thời gian ngắn nhất giữa hai lần tốc độ cực đại là một chu kì}
	{\True Độ lớn gia tốc cực đại trong quá trình dao động là $50 \ m/s^2$}
	\loigiai{
		\begin{itemchoice}
			\itemch Sai. Khi vật đi đến biên thì vận tốc vật bằng không.
			\itemch Sai. Tần số góc dao động là $\omega=\sqrt{\dfrac{k}{m}}=\sqrt{\dfrac{100}{0,1}}=10 \sqrt{10} rad/s$.
			\itemch Sai. Thời gian ngắn nhất giữa hai lần tốc độ cực đại là $1 / 2$ chu kì.
			\itemch Đúng.
		\end{\itemchoice}
	}
	%\haidong{25}
\end{ex}
\subsubsection*{PHẦN 3. Thí sinh trả lời từ câu 1 đến câu 6.}
\setcounter{ex}{0}
\begin{ex}
	Một vật dao động điều hoà trên trục $O x$ có vận tốc cực đại bằng $20 \ cm/s$ và gia tốc cực đại bằng $40\ cm/s^2$. Tốc độ góc của vật là bao nhiêu rad/s?
	\shortans[oly]{2}
	\loigiai{
		$2\ rad/s$
	}
	%\haidong{30}
\end{ex}
\begin{ex}
	Một vật dao động điều hòa với phương trình $x=8 \cos \left(2 \pi t-\dfrac{\pi}{6}\right)\ cm$. Thời điểm thứ $2015$ vật qua vị trí có vận tốc $-8 \pi\ cm/s$ là bao nhiêu phút (làm tròn kết quả đến chữ số hàng phần mười)?
	\shortans[oly]{16,8}
	\loigiai{
		16,8 phút
	}
	%\haidong{32}
\end{ex}
\begin{ex}
	Một vật nhỏ dao động điều hòa theo phương trình $x=A \cos 3 \pi t$ ($t$ tính bằng $s$). Thời điểm đầu tiên để gia tốc của vật có độ lớn bằng một nửa độ lớn gia tốc cực đại là bao nhiêu giây (làm tròn kết quả đến chữ số hàng phần trăm)?
	\shortans[oly]{0,11} \loigiai{
		0,11
	}
	%\haidong{30}
\end{ex}
\begin{ex}
	Động năng của một chất điểm dao động điều hoà được mô tả bởi phương trình sau: $W_{\text{đ}}=0,8 \sin ^2\left(6 \pi t+\dfrac{\pi}{6}\right)\ J$ (t tính bằng s). Thế năng của chất điểm tại thời điểm $t=1 \ s$ là bao nhiêu jun?
	\shortans[oly]{0,6}
	\loigiai{
		0,6
	}
	%\haidong{30}
\end{ex}
\begin{ex}
	Một chất điểm dao động điều hoà có hệ thức liên hệ giữa li độ $x$ và vận tốc $v$ là $\dfrac{x^2}{16}+\dfrac{v^2}{640}=1$ (trong đó $x$ tính bằng $cm$, $v$ tính bằng $cm/s)$. Lấy $\pi^2=10$. Chu kì dao động của chất điểm bằng bao nhiêu giây?
	\shortans[oly]{1}
	\loigiai{
		$1 \ s$
	}
	%\haidong{30}
\end{ex}
\begin{ex}%[3P1K6-2][Loai1]%Bài 15 dạng 1
	Một con lắc lò xo treo thẳng đứng gồm lò xo có độ cứng $100\ N/m$ và vật nhỏ có khối lượng $100\ g$. Giữ vật theo phương thẳng đứng làm lò xo giãn $3\ cm$ rồi truyền cho nó tốc độ $20\pi \sqrt{3}\ cm/s$ theo phương thẳng đứng, sau đó vật dao động điều hòa. Lấy $g = 10\ m/s^2,\ \pi^2=10$. Biên độ dao động của vật là bao nhiêu cm?
	\shortans[oly]{4}
	\loigiai{
		\begin{itemize}
			\item Ta có: $\omega  = 10\pi\  rad/s$ và $\Delta \ell_0 = 1\ cm$.
			\item Khi truyền tốc độ: lò xo đang giãn 3 cm $\Rightarrow$ vật ở dưới vị trí lò xo tự nhiên 3 cm $\Rightarrow$ vật đang cách vị trí cân bằng $\left|x\right| = \left|\Delta \ell_0 - 3\ cm\right| = 2\ cm$ và $\left|v\right| =20\pi \sqrt{3}\ cm/s \Rightarrow  A=\sqrt{x^2+\dfrac{v^2}{\omega^2}}=  4\ cm.$	
		\end{itemize}
	}
	%\haidong{30}
\end{ex}